\section{Olver's work}\label{Olver}
Olver \cite[(7.3)]{O} considers the dif\/ferential equation
\begin{gather}\label{1:eq1}
w''(z)=\frac1z w'(z)+\left(u^2+\frac{\mu^2-1}{z^2}+f(z)\right)w(z) .
\end{gather}
The function $f(z)$ is even and analytic in a simply-connected domain $D$ containing $0$.
It is assumed that $\Re\mu\ge 0$.
The goal is to f\/ind the asymptotic
behavior of solutions of \eqref{1:eq1} as $0<u\to\infty$.

Olver \cite[(7.4)]{O} starts with a formal solution to \eqref{1:eq1} of the form
\begin{gather*}%\label{1:eq2}
 w(z)=z\Z_\mu(uz)\sum_{s=0}^\infty \frac{A_s(z)}{u^{2s}}+\frac{z}{u}\Z_{\mu+1}(uz)\sum_{s=0}^\infty \frac{B_s(z)}{u^{2s}} ,
\end{gather*}
where either $\Z_\mu=I_\mu$, $\Z_{\mu+1}=I_{\mu+1}$ or $\Z_\mu=K_\mu$, $\Z_{\mu+1}=-K_{\mu+1}$ are modif\/ied Bessel functions.
The functions $A_s(z)=A_s(\mu,z)$, $B_s(z)=B_s(\mu,z)$ are def\/ined by $A_0(z)=1$, and then recursively, for $s\ge 0$,
\begin{gather}
2B_s(z) = -A_s'(z)+\int_0^z\left(f(t)A_s(t)-\frac{2\mu+1}{t} A_s'(t)\right) dt,\label{1:eq3}\\
2A_{s+1}(z) = \frac{2\mu+1}{z} B_s(z)- B_s'(z)+\int f(z)B_s(z) dz.\label{1:eq4}
\end{gather}
The integral in \eqref{1:eq4} denotes an arbitrary antiderivative of $f(z)B_s(z)$.
The func\-tions~$A_s(z)$,~$B_s(z)$ are analytic in $D$, and they are even and odd, respectively.

If the domain $D$ is unbounded, Olver \cite[p.~77]{O} requires that $f(z)=O(|z|^{-1-\alpha})$ as $|z|\to\infty$, where $\alpha>0$.
In our application to the conf\/luent hypergeometric equation in Section~\ref{confluent} the function $f(z)=z^2$ does not satisfy this condition.
Therefore, throughout this paper, we will take
\begin{gather}\label{1:D}
 D=\{z\colon |z|<R_0\},
 \end{gather}
where $R_0$ is a positive constant.
Olver \cite[p.~77]{O} introduces various subdomains $D'$, $D_1$, $D_2$ of~$D$.
We may choose $D'=\{z\colon |z|\le R\}$, where $0<R<R_0$. The domain $D_1$ comprises those points~$z$ in~$D'$ which can be joined to the origin by a contour
which lies in~$D'$ and does not cross either the imaginary axis, or the line through~$z$ parallel to the imaginary axis.
For our special~$D'$ the contour can be taken as the line segment connecting $z$ and $0$, so $D_1=D'$. The domain~$D_1$ appears in
Olver \cite[Theorem~D(i)]{O}. According to this theorem, \eqref{1:eq1} has a solu\-tion~$W_1(u,z)$ of the form
\begin{gather}\label{1:W1}
W_1(u,z)=zI_\mu(uz)\left(\sum_{s=0}^{N-1} \frac{A_s(z)}{u^{2s}}+g_1(u,z)\right)
+\frac{z}{u}I_{\mu+1}(uz)\left(\sum_{s=0}^{N-1} \frac{B_s(z)}{u^{2s}}+zh_1(u,z)\right),\!\!\!
\end{gather}
where
\begin{gather}\label{1:W1a}
 |g_1(u,z)|+|h_1(u,z)|\le \frac{K_1}{u^{2N}}\qquad \text{for} \quad 0<|z|\le R,\quad u\ge u_1 .
\end{gather}

\begin{Remarks}\quad
\begin{enumerate}\itemsep=0pt
\item
The parameter $\mu$ is considered f\/ixed. We may write $W_1(u,\mu,z)$ to indicate the dependence of $W_1$ on $\mu$.
\item
Every solution $w(z)$ of \eqref{1:eq1} is def\/ined on the Riemann surface of the logarithm over~$D$. Note that there is no restriction on $\arg z$ in~\eqref{1:W1a}, see \cite[p.~76]{O}.
\item
The precise statement is this: for every positive integer $N$ there are functions~$g_1$,~$h_1$ and positive constants~$K_1$, $u_1$ (independent of $u,z$)
such that~\eqref{1:W1},~\eqref{1:W1a} hold.
\item
The functions $A_s(z)$, $B_s(z)$ are not uniquely determined because of the free choice of
integration constants in~\eqref{1:eq4}. Even if we make a def\/inite choice of these integration constants, the solution $W_1(u,z)$
is not uniquely determined by~\eqref{1:W1},~\eqref{1:W1a}. For example, one can replace $W_1(u,z)$ by $(1+e^{-u})W_1(u,z)$.
\item
Olver's construction of $W_1(u,z)$ is independent of $N$ but may depend on~$R$.
In our application to the conf\/luent hypergeometric dif\/ferential equation we have $f(z)=z^2$. Then~$R$ can be any positive number
but $W_1(u,z)$ may depend on the choice of $R$.
\item
Olver has the term $\frac{z}{1+|z|}$ in place of~$z$ in front of~$h_1$ in~\eqref{1:W1} but since we assume $|z|\le R$ this makes no dif\/ference.
\end{enumerate}
\end{Remarks}

For the def\/inition of $D_2$ we suppose that $a$ is an arbitrary point of the sector $|\arg a|<\frac12\pi$ and $\epsilon>0$.
Then $D_2$ comprises those points $z\in D'$ for which $|\arg z|\le\frac32\pi$, $\Re z\le \Re a$, and a~contour can be found joining~$z$ and~$a$
which satisf\/ies the following conditions:
\begin{itemize}\itemsep=0pt
\item[(i)] it lies in $D'$,
\item[(ii)]
it lies wholly to the right of the line through $z$ parallel to the imaginary axis,
\item[(iii)]
it does not cross the negative imaginary axis if $\frac12\pi\le \arg z\le \frac32\pi$, and does not cross
the positive imaginary axis if $-\frac32\pi\le \arg z\le -\frac\pi2$,
\item[(iv)]
it lies outside the circle $|t|=\epsilon |z|$.
\end{itemize}

In our special case $D'=\{z\colon |z|\le R\}$ we choose $a=R$. If $0\le\arg z\le \frac32\pi$ and $0<|z|\le R$, we choose the contour starting at~$z$
moving in positive direction parallel to the imaginary axis until we hit the circle $|t|=R$. Then we move clockwise along the circle $|t|=R$ towards~$a$. Taking into account condition~(iv), we see that~$D_2$ is the set of points $z$ with $-\frac32\pi+\delta\le z\le \frac32\pi-\delta$,
$0<|z|\le R$, where $\delta>0$.
The domain $D_2$ appears in Olver \cite[Theorem~D(ii)]{O}. According to this theorem, \eqref{1:eq1} has a~solution $W_2(u,z)$ of the form
\begin{gather}
W_2(u,z)=zK_\mu(uz)\left(\sum_{s=0}^{N-1} \frac{A_s(z)}{u^{2s}}+g_2(u,z)\right)\nonumber\\
\hphantom{W_2(u,z)=}{} -\frac{z}{u}K_{\mu+1}(uz)\left(\sum_{s=0}^{N-1} \frac{B_s(z)}{u^{2s}}+zh_2(u,z)\right),\label{1:W2}
\end{gather}
where
\begin{gather}\label{1:W2a}
 |g_2(u,z)|+|h_2(u,z)|\le \frac{K_2}{u^{2N}}\qquad \text{for} \quad 0<|z|\le R, \quad |\arg z|\le \frac32\pi-\delta, \quad u\ge u_2 .
\end{gather}
Note that in \eqref{1:W2a} there is a restriction on $\arg z$.

In the rest of this paper we choose the functions $A_s(z)$ such that
\begin{gather}\label{1:As}
A_s(0)=0\qquad\text{if} \quad s\ge 1.
\end{gather}
Then the functions $A_s(z)$, $B_s(z)$ are uniquely determined.

\section[Properties of solutions $W_1$ and $W_2$]{Properties of solutions $\boldsymbol{W_1}$ and $\boldsymbol{W_2}$}\label{discussion}

The dif\/ferential equation \eqref{1:eq1} has a regular singularity at $z=0$ with exponents $1\pm \mu$.
Substituting $x=z^2$ we obtain an equation which has a regular singularity at $x=0$ with exponents $\frac12(1\pm\mu)$.
Therefore, for every $\mu$ which is not a negative integer, \eqref{1:eq1} has a unique solution $W_+(z)=W_+(u,\mu,z)$ of the form
\begin{gather*}%\label{2:eq1}
W_+(z)= z^{1+\mu}\sum_{n=0}^\infty c_nz^{2n},
\end{gather*}
where the $c_n$ are determined by $c_0=1$, and
\begin{gather*}
4n(\mu+n)c_n=u^2c_{n-1}+\sum_{j=0}^{n-1}f_jc_{n-1-j} \qquad\text{for} \quad n\ge 1
\end{gather*}
when
\begin{gather*}
f(z)=\sum_{n=0}^\infty f_n z^{2n} .
\end{gather*}

If $\mu$ is not an integer, then $W_+(u,\mu,z)$ and $W_+(u,-\mu,z)$ form a fundamental system of solutions of~\eqref{1:eq1}.
If $\Re \mu\ge 0$, there is a solution $W_-(z)$ linearly independent of $W_+(z)$ such that
\begin{gather*}%\label{2:eq2}
W_-(z)=z^{1-\mu}p\big(z^2\big)+ d \ln z W_+(z),
\end{gather*}
where $p$ is a power series and $d$ is a suitable constant. If $\mu\ne 0$ we choose $p(0)=1$. If~$\mu$ is not an integer then $d=0$.

\begin{Lemma}\label{2:l1}
Suppose $\Re\mu\ge 0$.
There is a function $\alpha(u)$ such that
\begin{gather*}%\label{2:connectW1}
W_1(u,z)=\alpha(u)W_+(u,z),
\end{gather*}
and, for every $N=1,2,3,\dots$,
\begin{gather}\label{2:alpha}
\alpha(u)=\frac{2^{-\mu}u^\mu}{\Gamma(\mu+1)}\left(1+O\left(\frac1{u^{2N}}\right)\right)\qquad\text{as} \quad 0<u\to\infty.
\end{gather}
\end{Lemma}

\begin{proof}
There are functions $\alpha_+(u)$, $\alpha_-(u)$ such that
\begin{gather}\label{2:connectW1a}
W_1(u,z)=\alpha_+(u)W_+(u,z)+\alpha_-(u)W_-(u,z).
\end{gather}
Suppose $\Re\mu>0$. Then \eqref{2:connectW1a} implies
\begin{gather}\label{2:limit1}
\lim_{z\to0^+} z^{\mu-1}W_1(u,z)=\alpha_-(u).
\end{gather}
We use \cite[(10.30.1)]{NIST}
\begin{gather*}%\label{2:limitI}
 \lim_{z\to0} I_\nu(z)z^{-\nu}=\frac{2^{-\nu}}{\Gamma(\nu+1)}.
\end{gather*}
Then \eqref{1:W1}, \eqref{1:W1a} give{\samepage
\begin{gather}\label{2:limit2}
\lim_{z\to0^+} z^{\mu-1}W_1(u,z)= 0.
\end{gather}
It follows from \eqref{2:limit1}, \eqref{2:limit2} that $\alpha_-(u)=0$.}

